% !TEX program = xelatex
%% ============================================================
%%  IL SALOTTO DI ENEA ED ETTORE – Piccolo manuale di carte per una notte ruggente
%%  ESN Modena e Reggio Emilia · 20 anni · Gala Anni '20
%%  Compilare con XeLaTeX (2 passate):   xelatex main.tex
%%  Formato: A5, 24 pagine. Per la stampa a libretto vedi LEGGIMI.txt.
%% ============================================================
\documentclass[10pt]{article}
%% ============================================================
%%  STILE GRAFICO – «Il salotto di Enea ed Ettore»
%%  Tutto ciò che riguarda aspetto, colori, font, carte e riquadri.
%%  Per cambiare la grafica basta toccare questo file.
%% ============================================================

% ---------- Pagina ----------
\usepackage{fix-cm}
\usepackage[paperwidth=148mm,paperheight=210mm,
            left=14mm,right=14mm,top=13mm,bottom=18mm]{geometry}
\usepackage{polyglossia}
\setdefaultlanguage{italian}
\usepackage{fontspec}
\usepackage{xcolor,graphicx,eso-pic,enumitem,colortbl,tabularx,array}
\usepackage[most]{tcolorbox}
\usepackage{tikz}
\usetikzlibrary{calc,shapes.geometric}
\usepackage{microtype}
\frenchspacing
\setlength{\parindent}{0pt}
\setlength{\parskip}{1.2mm}
\setlength{\emergencystretch}{3em}
\pagestyle{empty}

% ---------- Colori (presi dalla mail e dalle grafiche ESN MoRe) ----------
\definecolor{ink}{HTML}{0D0D0B}
\definecolor{gold}{HTML}{C9953F}
\definecolor{golddk}{HTML}{9A6B1C}
\definecolor{cream}{HTML}{EFE3CC}
\definecolor{paper}{HTML}{F7F0DF}
\definecolor{tint}{HTML}{EFE0BD}
\definecolor{rosepig}{HTML}{E8A9B5}
\definecolor{rose}{HTML}{B8607A}
\definecolor{bluepig}{HTML}{8DB3C6}
\definecolor{teal}{HTML}{10262B}
\definecolor{cardred}{HTML}{B3202A}
\definecolor{txt}{HTML}{2A2418}

% ---------- Font (nella cartella fonts/) ----------
\setmainfont{JosefinSans}[Path=./fonts/,Extension=.ttf,
  UprightFont=*-Regular,BoldFont=*-Bold,ItalicFont=*-Italic,BoldItalicFont=*-BoldItalic]
\newfontfamily\deco{PoiretOne-Regular.ttf}[Path=./fonts/]
\newfontfamily\decob{PoiretOne-Regular.ttf}[Path=./fonts/,AutoFakeBold=3,
  ItalicFont=PoiretOne-Regular.ttf,BoldItalicFont=PoiretOne-Regular.ttf]
\newfontfamily\suitfont{DejaVuSans.ttf}[Path=./fonts/,BoldFont=DejaVuSans-Bold.ttf]
\newcommand{\corpo}{\fontsize{8.9}{12.2}\selectfont\color{txt}}
% Celle ed esempi: dimensioni condivise, contenuto centrato su entrambi gli assi.
\newcolumntype{C}[1]{>{\centering\arraybackslash}m{#1}}
\newcommand{\examplecell}[3]{\begin{minipage}[t]{#1}\centering
  \vspace{0pt}\makebox[\linewidth][c]{\raisebox{0pt}[19mm][0pt]{#2}}%
  \par\vspace{1.2mm}{\footnotesize #3}\end{minipage}}
\newcommand{\cardref}[2]{\mbox{#1\,#2}}
\newcommand{\headerdiamond}{\tikz[baseline=-.05ex]{\fill[gold]
  (0,.5ex)--(.5ex,1ex)--(1ex,.5ex)--(.5ex,0)--cycle;}}

% Pagine da compilare a mano: la stessa griglia per quattro o sei giocatori.
\newcommand{\scoresheet}[2]{%
 \startpage\label{segnapunti-#1}
 \bigtitle{Segnapunti}{#2}
 \vspace{3mm}
 Gioco: \rule{48mm}{.3pt}\hfill Tavolo: \rule{29mm}{.3pt}\par
 \vspace{3mm}
 {\footnotesize Scrivete i nomi e segnate il risultato di ogni mano.}
 \begin{center}
 \begin{tikzpicture}[x=1mm,y=1mm,font=\small]
  \fill[tint] (0,0) rectangle (114,-8);
  \fill[tint] (0,-104) rectangle (114,-112);
  \foreach \y in {0,-8,...,-112}{
   \draw[gold,line width=.4pt] (0,\y)--(114,\y);}
  \draw[gold,line width=.4pt] (0,0)--(0,-112) (18,0)--(18,-112);
  \foreach \i in {1,...,#1}{
   \draw[gold,line width=.4pt] ({18+96*\i/#1},0)--({18+96*\i/#1},-112);
   \node at ({18+96*(\i-.5)/#1},-4) {Nome};}
  \node[font=\small\bfseries] at (9,-4) {Mano};
  \foreach \i in {1,...,12}{\node at (9,{-4-8*\i}) {\i};}
  \node[font=\small\bfseries] at (9,-108) {Totale};
 \end{tikzpicture}
 \end{center}
 \vfill
 \begin{center}\itshape\small\color{golddk}
 Una partita in più, un’altra storia da raccontare.
 \end{center}}

\newcommand{\writinglines}[2]{%
 \par\noindent\begin{tikzpicture}[x=1mm,y=1mm]
  \useasboundingbox (0,0) rectangle (120,{-#1*#2});
  \foreach \k in {1,...,#1}{%
   \draw[gold!55!paper,line width=.35pt] (.2,{-\k*#2})--(119.8,{-\k*#2});}
 \end{tikzpicture}\par}

% ---------- Sfondi di pagina ----------
\gdef\pagetype{cream}
\newcommand{\startpage}[1][cream]{\clearpage\gdef\pagetype{#1}}

\newcommand{\decoframe}[4]{% inset, smusso, colore, spessore
 \draw[#3,line width=#4] (#1+#2,#1) -- (\paperwidth-#1-#2,#1) -- (\paperwidth-#1,#1+#2)
   -- (\paperwidth-#1,\paperheight-#1-#2) -- (\paperwidth-#1-#2,\paperheight-#1)
   -- (#1+#2,\paperheight-#1) -- (#1,\paperheight-#1-#2) -- (#1,#1+#2) -- cycle;}
\newcommand{\cornerfan}[4]{% x,y,rotazione,colore
 \begin{scope}[shift={(#1,#2)},rotate=#3]
  \foreach \a in {0,15,...,90}{\draw[#4,line width=.3pt] (0,0) -- (\a:8mm);}
  \draw[#4,line width=.3pt] (8mm,0) arc(0:90:8mm);
  \draw[#4,line width=.3pt] (5mm,0) arc(0:90:5mm);
 \end{scope}}
\newcommand{\framecorners}[2]{%
 \cornerfan{#1}{#1}{0}{#2}
 \cornerfan{\paperwidth-#1}{#1}{90}{#2}
 \cornerfan{\paperwidth-#1}{\paperheight-#1}{180}{#2}
 \cornerfan{#1}{\paperheight-#1}{270}{#2}}

\newcommand{\bgcream}{%
 \begin{tikzpicture}
  \useasboundingbox (0,0) rectangle (\paperwidth,\paperheight);
  \fill[paper] (0,0) rectangle (\paperwidth,\paperheight);
  \decoframe{6mm}{3.5mm}{gold}{.7pt}
  \decoframe{7.6mm}{3.5mm}{gold}{.3pt}
  \framecorners{6mm}{gold}
  % numero di pagina in un rombo di misura fissa (una o due cifre):
  % la punta inferiore resta sopra il filetto interno della cornice (7,6 mm)
  \node[diamond,draw=gold,fill=ink,line width=.5pt,inner sep=0pt,minimum size=6.4mm,
        text=cream,font=\small,text height=1.5ex,text depth=.2ex] at (\paperwidth/2,11.6mm) {\thepage};
  \draw[gold,line width=.4pt] (\paperwidth/2-18mm,11.6mm) -- (\paperwidth/2-5.5mm,11.6mm);
  \draw[gold,line width=.4pt] (\paperwidth/2+5.5mm,11.6mm) -- (\paperwidth/2+18mm,11.6mm);
 \end{tikzpicture}}
\newcommand{\bgdark}{%
 \begin{tikzpicture}
  \useasboundingbox (0,0) rectangle (\paperwidth,\paperheight);
  \fill[black] (0,0) rectangle (\paperwidth,\paperheight);
  \decoframe{6mm}{3.5mm}{gold}{.8pt}
  \decoframe{7.8mm}{3.5mm}{gold}{.3pt}
  \framecorners{6mm}{gold}
 \end{tikzpicture}}

% Cornice che sfuma verso il basso (solo copertina).
% Argomenti: inset, smusso, colore, spessore, y_pieno, y_zero
% y_pieno e y_zero sono in mm dal BORDO INFERIORE, senza unità:
% sopra y_pieno la cornice è piena, sotto sfuma fino a sparire a y_zero.
% Misure fisse per il formato A5 (148 x 210 mm).
\newcommand{\fadeframe}[6]{%
	\pgfmathsetmacro{\xl}{#1}\pgfmathsetmacro{\xr}{148-#1}%
	\pgfmathsetmacro{\xls}{#1+#2}\pgfmathsetmacro{\xrs}{148-#1-#2}%
	\pgfmathsetmacro{\yt}{210-#1}\pgfmathsetmacro{\ytop}{210-#1-#2}%
	\pgfmathsetmacro{\ystep}{(210-2*#1-2*#2)/120}%
	\draw[#3,line width=#4] (\xl mm,\ytop mm) -- (\xls mm,\yt mm) -- (\xrs mm,\yt mm) -- (\xr mm,\ytop mm);%
	\foreach \i in {0,...,119}{%
		\pgfmathsetmacro{\ya}{\ytop-\ystep*\i}%
		\pgfmathsetmacro{\yb}{\ytop-\ystep*(\i+1)}%
		\pgfmathsetmacro{\op}{max(0,min(1,((\ya+\yb)/2-#6)/(#5-#6)))}%
		\draw[#3,line width=#4,opacity=\op] (\xl mm,\ya mm) -- (\xl mm,\yb mm);%
		\draw[#3,line width=#4,opacity=\op] (\xr mm,\ya mm) -- (\xr mm,\yb mm);%
}}

% L'immagine è un nodo dentro la stessa tikzpicture: prima la cornice
% veniva disegnata dopo l'immagine, cioè fuori dal foglio, e non si vedeva.
\newcommand{\bgcover}{%
 \begin{tikzpicture}
  \useasboundingbox (0,0) rectangle (\paperwidth,\paperheight);
  \node[anchor=south west,inner sep=0pt,outer sep=0pt] at (0,0)
    {\includegraphics[width=\paperwidth,height=\paperheight]{img/cover_bg.png}};
	\fadeframe{6}{3.5}{gold}{.8pt}{175}{120}
	\fadeframe{7.8}{3.5}{gold}{.3pt}{175}{120}
	\cornerfan{6mm}{\paperheight-6mm}{270}{gold}
	\cornerfan{\paperwidth-6mm}{\paperheight-6mm}{180}{gold}
 \end{tikzpicture}}
\AddToShipoutPictureBG{\AtPageLowerLeft{%
  \csname bg\pagetype\endcsname}}

% ---------- Intestazione di gioco (fascia scura a tutta pagina) ----------
\newcommand{\diffd}[1]{\tikz[baseline=-.05ex]{%
  \foreach \k in {1,...,5}{\ifnum\k>#1\def\dfill{none}\else\def\dfill{gold}\fi
    \path[draw=gold,fill=\dfill,line width=.35pt]
      ({(\k-1)*1.8ex},.7ex)--++(.7ex,.7ex)--++(.7ex,-.7ex)--++(-.7ex,-.7ex)--cycle;}}}
% Centra ogni elemento sulla sua altezza effettiva, indipendentemente dal font.
\newcommand{\headmid}[2][0pt]{\begingroup
 \setbox0=\hbox{#2}%
 \raisebox{\dimexpr(\dp0-\ht0)/2+#1\relax}{\box0}%
 \endgroup}
\newcommand{\gameheader}[7]{% num, TITOLO, sottotitolo, giocatori, durata, difficoltà, mascotte
\begin{tikzpicture}[remember picture,overlay]
 \fill[ink] (current page.north west) rectangle ([yshift=-47mm]current page.north east);
 \draw[gold,line width=.7pt] ([yshift=-44.5mm]current page.north west) -- ([yshift=-44.5mm]current page.north east);
 \draw[gold,line width=.3pt] ([yshift=-46mm]current page.north west) -- ([yshift=-46mm]current page.north east);
 \begin{scope}[shift={(current page.north west)}]
  \cornerfan{0}{0}{270}{gold!55!ink}
 \end{scope}
 \begin{scope}[shift={(current page.north east)}]
  \cornerfan{0}{0}{180}{gold!55!ink}
 \end{scope}
 \node[anchor=base west,text=gold,font=\decob\bfseries\footnotesize,inner sep=0] at ([xshift=14mm,yshift=-8.5mm]current page.north west)
   {\headmid{\ifstrequal{#1}{BONUS}{BONUS}{N°~#1}}%
    \hspace{1.5mm}\headmid[.32mm]{\headerdiamond}\hspace{1.5mm}\headmid{\MakeUppercase{#4}}};
 \node[anchor=north west,text=gold,font=\decob\bfseries\fontsize{30}{32}\selectfont,inner sep=0] at ([xshift=14mm,yshift=-11mm]current page.north west)
   {#2};
 \node[anchor=north west,text=cream,font=\itshape\small,inner sep=0] at ([xshift=14mm,yshift=-22mm]current page.north west)
   {#3};
 % Centro verticale comune, con correzioni ottiche per l’inchiostro visibile
 % di Poiret/Josefin e dei simboli (verificate sul PDF renderizzato).
 \node[anchor=base west,text=gold,font=\decob\bfseries\footnotesize,inner sep=0]
   at ([xshift=14mm,yshift=-31mm]current page.north west) {\headmid{DURATA}};
 \node[anchor=base west,text=cream,font=\footnotesize,inner sep=0]
   at ([xshift=32mm,yshift=-31mm]current page.north west) {\headmid[.16mm]{#5}};
 \node[anchor=base west,text=gold,font=\decob\bfseries\footnotesize,inner sep=0]
   at ([xshift=14mm,yshift=-35.5mm]current page.north west) {\headmid{DIFFICOLTÀ}};
 \node[anchor=base west,text=gold,font=\footnotesize,inner sep=0]
   at ([xshift=32mm,yshift=-35.5mm]current page.north west) {\headmid[.53mm]{\diffd{#6}}};
 \node[circle,fill=ink,inner sep=0] at ([xshift=-26mm,yshift=-33mm]current page.north east)
   {\includegraphics[width=32mm]{img/#7.png}};
\end{tikzpicture}\vspace*{37mm}\par\nopagebreak}

% ---------- Titoli di sezione ----------
\newcommand{\secz}[1]{\par\nobreak\vspace{1.2mm}%
 {\noindent\decob\bfseries\fontsize{12.5}{14}\selectfont\color{ink}\MakeUppercase{#1}\hspace{2mm}%
  {\color{gold}\leaders\hrule height .5pt\hfill}\kern0pt\par}\nobreak\vspace{.4mm}}
\newcommand{\bigtitle}[2]{%
 \begin{center}{\decob\bfseries\fontsize{23}{26}\selectfont\color{ink}\MakeUppercase{#1}}\\[1mm]
 {\itshape\small\color{golddk}#2}\end{center}}
\newcommand{\divider}{\par\vspace{1mm}\noindent\begin{center}\begin{tikzpicture}
  \draw[gold,line width=.4pt] (-24mm,0) -- (-3mm,0); \draw[gold,line width=.4pt] (3mm,0) -- (24mm,0);
  \node[diamond,fill=gold,inner sep=1.2pt] at (0,0) {};
 \end{tikzpicture}\end{center}\par\vspace{.5mm}}

% ---------- Elenchi ----------
\setlist{nosep}
\setlist[enumerate]{label=\textcolor{golddk}{\decob\bfseries\arabic*}\,,leftmargin=5.2mm,itemsep=1mm,topsep=.8mm,parsep=0pt}
\setlist[itemize]{label=\raisebox{.2ex}{\textcolor{gold}{\suitfont\scriptsize ◆}},leftmargin=4.6mm,itemsep=.7mm,topsep=.8mm,parsep=0pt}
\setlist[itemize,2]{label=\textcolor{gold}{–},leftmargin=4mm,itemsep=.3mm,topsep=.3mm}

% ---------- Riquadri ----------
\newtcolorbox{callout}[1][]{enhanced,breakable=false,colback=tint,colframe=tint,boxrule=0pt,arc=0pt,
  left=3.6mm,right=2.6mm,top=1.5mm,bottom=1.5mm,borderline west={1.8pt}{0pt}{gold},
  before skip=2mm,after skip=2mm,fontupper=\footnotesize\color{txt},#1}
\newcommand{\tipbox}[2]{% chi = enea | ettore
 \begin{tcolorbox}[enhanced,breakable=false,colback=ink,colframe=gold,boxrule=.6pt,arc=2mm,
   left=2mm,right=2.4mm,top=1.2mm,bottom=1.2mm,before skip=1.6mm,after skip=1mm,coltext=cream]
 \begin{minipage}[c]{12mm}\includegraphics[width=11.5mm]{img/#1.png}\end{minipage}\hfill
 \begin{minipage}[c]{\dimexpr\linewidth-14.5mm\relax}\footnotesize #2\end{minipage}
 \end{tcolorbox}}
\newcommand{\tipname}[1]{\textcolor{gold}{\decob\bfseries\large #1}\\[.4mm]}
\newcommand{\enea}[1]{\tipbox{enea}{\tipname{\textcolor{rosepig}{Il consiglio di Enea}}#1}}
\newcommand{\ettore}[1]{\tipbox{ettore}{\tipname{\textcolor{bluepig}{Il trucco di Ettore}}#1}}
\newcommand{\termine}[1]{\textbf{#1}}

% ---------- Carte da gioco in TikZ ----------
\makeatletter
\def\suit@S{♠}\def\suit@C{♣}\def\suit@H{♥}\def\suit@D{♦}\def\suit@Z{★}
\def\scol@S{ink}\def\scol@C{ink}\def\scol@H{cardred}\def\scol@D{cardred}\def\scol@Z{rose}
\def\x@back{X}
\newcommand{\cw}{11mm}\newcommand{\cpitch}{9.2mm}
% \cardnode{valore}{seme}{x}{modo}  modo: n normale, u evidenziata (sollevata), d scartata (sbiadita)
\newcommand{\cardnode}[4]{%
	\edef\c@rank{#1}\expandafter\let\expandafter\c@suit\csname suit@#2\endcsname
	\expandafter\let\expandafter\c@col\csname scol@#2\endcsname
	\edef\c@mode{#4}\def\c@u{u}\def\c@d{d}%
	\begin{scope}[shift={(#3,0)}]
		\ifx\c@mode\c@u\begin{scope}[yshift=2.6mm]\fi
			\ifx\c@rank\x@back
			\draw[fill=ink,draw=gold,line width=.5pt,rounded corners=.9mm] (0,0) rectangle (\cw,{1.4*\cw});
			\draw[gold,line width=.3pt,rounded corners=.4mm] (1.2mm,1.2mm) rectangle ({\cw-1.2mm},{1.4*\cw-1.2mm});
			\foreach \a in {20,50,...,160}{\draw[gold,line width=.3pt] ({.5*\cw},{.45*\cw}) -- +(\a:6mm);}
			\node[diamond,fill=gold,inner sep=1.4pt] at ({.5*\cw},{.7*\cw}) {};
			\else
			\draw[fill=white,draw=ink!65,line width=.3pt,rounded corners=.9mm] (0,0) rectangle (\cw,{1.4*\cw});
			\ifx\c@u\c@mode\draw[gold,line width=1.1pt,rounded corners=.9mm] (0,0) rectangle (\cw,{1.4*\cw});\fi
			\def\c@z{Z}\edef\c@s{#2}
			\ifx\c@s\c@z
			\node at ({.5*\cw},{.7*\cw}) {\suitfont\fontsize{20}{20}\selectfont\color{\c@col}\c@suit};
			\node[anchor=north west,inner sep=0] at (1mm,{1.4*\cw-.5mm}) {\fontsize{5.2}{5.2}\selectfont\bfseries\color{rose}JOLLY};
			\else
			% --- Angolo in alto a sinistra ---
			\node[anchor=north,inner sep=0] at (1.9mm,{1.4*\cw-.6mm}) 
			{\fontsize{9}{9}\selectfont\bfseries\color{\c@col}\c@rank};
			\node[anchor=north,inner sep=0] at (1.9mm,{1.4*\cw-3.5mm}) 
			{\suitfont\fontsize{6.5}{6.5}\selectfont\color{\c@col}\c@suit};
			
			% --- Angolo in basso a destra ---
			\node[anchor=north,inner sep=0,rotate=180] at ({\cw-1.9mm},.6mm) 
			{\fontsize{9}{9}\selectfont\bfseries\color{\c@col}\c@rank};
			\node[anchor=north,inner sep=0,rotate=180] at ({\cw-1.9mm},{3.5mm}) 
			{\suitfont\fontsize{6.5}{6.5}\selectfont\color{\c@col}\c@suit};
			
			% --- Simbolo o Figura centrale (pulita) ---
			\def\c@J{J}\def\c@Q{Q}\def\c@K{K}
			\ifx\c@rank\c@J\def\c@face{1}\else\ifx\c@rank\c@Q\def\c@face{1}\else\ifx\c@rank\c@K\def\c@face{1}\else\def\c@face{0}\fi\fi\fi
			\if\c@face1
			\node at ({.5*\cw},{.7*\cw}) {\bfseries\fontsize{15}{15}\selectfont\color{\c@col}\c@rank};
			\else
			\node at ({.5*\cw},{.7*\cw}) {\suitfont\fontsize{16}{16}\selectfont\color{\c@col}\c@suit};
			\fi
			\fi
			\fi
			\ifx\c@mode\c@d\fill[cream,opacity=.68,rounded corners=.9mm] (0,0) rectangle (\cw,{1.4*\cw});\fi
			\ifx\c@mode\c@u\end{scope}\fi
\end{scope}}
\makeatother
% \hand{V/S,V/S,...}  oppure  \handm{V/S/m,...}
\newcommand{\hand}[1]{\begin{tikzpicture}[baseline=0pt]
 \foreach \r/\s [count=\i from 0] in {#1}{\cardnode{\r}{\s}{\i*\cpitch}{n}}
 \end{tikzpicture}}
\newcommand{\handm}[1]{\begin{tikzpicture}[baseline=0pt]
 \foreach \r/\s/\m [count=\i from 0] in {#1}{\cardnode{\r}{\s}{\i*\cpitch}{\m}}
 \end{tikzpicture}}
% versione ridotta: \hands{scala}{lista}
\newcommand{\handsc}[2]{\scalebox{#1}{\hand{#2}}}
\newcommand{\sS}{{\suitfont ♠}}\newcommand{\sC}{{\suitfont ♣}}
\newcommand{\sH}{{\suitfont\color{cardred}♥}}\newcommand{\sD}{{\suitfont\color{cardred}♦}}
\newcommand{\fiche}{\tikz[baseline=-.6ex]{\draw[fill=gold,draw=ink,line width=.4pt] (0,0) circle (1.6mm);\draw[ink,line width=.3pt,dash pattern=on .6mm off .5mm] (0,0) circle (1.05mm);}}
\newcommand{\ruleh}{\par\vspace{.5mm}}

\begin{document}
\corpo

%% ------------------------------------------------------------
%% 1 · COPERTINA
%% ------------------------------------------------------------
\gdef\pagetype{cover}
\thispagestyle{empty}
\begin{tikzpicture}[remember picture,overlay]
	\node[text=gold,font=\deco\fontsize{8.5}{10}\selectfont,align=center] at ([yshift=-15mm]current page.north)
	{20 ANNI DI ESN MORE\quad\headmid[.32mm]{\headerdiamond}\quad GALA ANNI ’20};
	\draw[gold,line width=.4pt] ([xshift=-30mm,yshift=-20mm]current page.north) -- ([xshift=-3mm,yshift=-20mm]current page.north);
	\draw[gold,line width=.4pt] ([xshift=3mm,yshift=-20mm]current page.north) -- ([xshift=30mm,yshift=-20mm]current page.north);
	\node[diamond,fill=gold,inner sep=1.3pt] at ([yshift=-20mm]current page.north) {};
	\node[text=gold,font=\decob\bfseries\fontsize{50}{52}\selectfont] at ([yshift=-35mm]current page.north) {IL SALOTTO DI};
	\node[text=gold,font=\decob\bfseries\fontsize{38}{40}\selectfont] at ([yshift=-51mm]current page.north) {ENEA ED ETTORE};
	\node[text=cream,font=\itshape\fontsize{10.5}{13}\selectfont,align=center] at ([yshift=-64mm]current page.north)
	{Il piccolo manuale di carte\\per una notte ruggente};
	\node[text=gold,font=\deco\fontsize{7.5}{9}\selectfont,align=center,draw=gold,line width=.4pt,
	      inner xsep=2.5mm,inner ysep=1.4mm,rounded corners=1mm]
	at ([yshift=-76mm]current page.north)
	{SEI GIOCHI\quad\headmid[.32mm]{\headerdiamond}\quad NESSUNA ESPERIENZA RICHIESTA\quad\headmid[.32mm]{\headerdiamond}\quad 21 NOVEMBRE 2026};
\end{tikzpicture}


%% 2 · BENVENUTI
\startpage
\begin{center}\vspace*{-1mm}\includegraphics[width=34mm]{img/duo.png}\end{center}
\vspace{-3mm}
\bigtitle{Benvenuti nel salotto}{Enea ed Ettore vi fanno gli onori di casa}
Gli anni Venti sono stati il decennio delle luci soffuse, dei velluti, delle piume e del jazz — e delle lunghe partite a carte, nei salotti dell’alta società come negli speakeasy.

Stasera tocca a voi. Il mazzo che avete tra le mani è il vostro biglietto d’ingresso, e questo libretto è la vostra guida: \textbf{sei giochi} spiegati passo dopo passo, dal più semplice al più astuto. Non serve esperienza: bastano un tavolo, qualche amico e un pizzico di sfrontatezza.

\secz{Come si usa}
\begin{itemize}
	\item I giochi sono in \textbf{ordine di difficoltà}: cominciate dal primo e salite quando vi sentite pronti.
	\item Ogni gioco ha la sua scheda: durata e difficoltà ({\suitfont\color{gold}◆} = facilissimo, 5 = da campioni).
	\item Nei riquadri scuri, \textbf{Enea} ed \textbf{Ettore} svelano consigli e trucchi.
	\item Si gioca \textbf{solo per gloria}: niente soldi, solo fiches. 
\end{itemize}
\begin{callout}
	\textbf{Le regole della serata.} Qui trovate versioni semplificate: alcune regole possono differire da quelle che conoscete. Seguite la stessa scheda per tutta la partita. 
\end{callout}

\secz{Cosa vi serve}
\begin{itemize}
	\item Il mazzo da \textbf{52 carte} (togliete i jolly: servono solo nel gioco bonus).
	\item Una ventina di \textbf{fiches} a testa.
	\item Carta e penna per i punti (i segnapunti sono alle pagine 20–22; segue una pagina per gli appunti).
	\item Per il gioco \textbf{bonus}: un secondo mazzo. Chiedetelo a un altro tavolo e fate amicizia!
\end{itemize}
\vspace{1mm}
\begin{center}
	\itshape\small\color{golddk}Mettete la maschera e lucidate il bastone da passeggio: si gioca.
\end{center}

%% 3 · PROGRAMMA
\startpage
\bigtitle{Il programma della serata}{cinque giochi in crescendo, più un bonus}
\divider
\begin{center}
\begin{tikzpicture}[x=1mm,y=1mm]
\draw[gold,line width=.5pt] (.2,4)--(119.8,4) (.2,-4)--(119.8,-4);
\foreach \x/\t in {3/N°,33.9/GIOCO,70.8/GIOCATORI,93.7/DIFFICOLTÀ,113.8/PAG.}{
 \node[text=golddk,font=\deco\footnotesize] at (\x,0) {\t};}
\newcommand{\indexrow}[7]{%
 \node[font=\small] at (3,{2-13*#1}) {#2};
 \node[align=center,text width=51mm,font=\small] at (33.9,{2-13*#1})
   {\textbf{#3}\\[.8mm]{\footnotesize\itshape #4}};
 \node[font=\small] at (70.8,{2-13*#1}) {#5};
 \node at (93.7,{2-13*#1}) {\diffd{#6}};
 \node[font=\small] at (113.8,{2-13*#1}) {\pageref{#7}};
 \draw[gold,line width=.35pt] (.2,{-4.5-13*#1})--(119.8,{-4.5-13*#1});}
\indexrow{1}{1}{Acey Deucey}{Tre carte, una scommessa}{3–8}{1}{acey}
\indexrow{2}{2}{Ventuno (Blackjack)}{Battere il banco senza sballare}{2–8}{1}{ventuno}
\indexrow{3}{3}{Poker a 5 carte}{Il brivido del bluff}{2–6}{2}{poker}
\indexrow{4}{4}{Gin Rummy}{Il duello delle combinazioni}{2}{3}{gin}
\indexrow{5}{5}{Cuori (Hearts)}{La Donna di Picche da evitare}{4 (o 3)}{3}{cuori}
\indexrow{6}{{\suitfont\color{gold}★}}{BONUS · Ramino}{Due mazzi, un tavolo di amici}{2–5}{2}{ramino}
\end{tikzpicture}
\end{center}
\vspace{3mm}
\secz{Da dove cominciare?}
\begin{itemize}
 \item \textbf{Siete in tanti e volete ridere:} Acey Deucey o Ventuno.
 \item \textbf{Volete il brivido del bluff:} Poker.
 \item \textbf{Un duello a due:} Gin Rummy.
 \item \textbf{Quattro amici e una serata lunga:} Cuori.
\end{itemize}
\vfill
\enea{Date prima un’occhiata alla preparazione: così scegliete il gioco giusto per il vostro tavolo e avete tutto a portata di mano.}

%% 4 · IL MAZZO IN BREVE
\startpage
\bigtitle{Il mazzo in breve}{tutto quello che serve sapere prima di dare le carte}
\divider
\begin{minipage}[c]{34mm}\centering
 \scalebox{1.8}{\hand{Q/H}}\par\vspace{1.5mm}
 {\footnotesize Q = Donna\quad\sH{} = Cuori}
\end{minipage}\hfill
\begin{minipage}[c]{82mm}\raggedright
 Il mazzo ha \textbf{52 carte}: quattro \textbf{semi} da 13 carte ciascuno.\par\medskip
 \sH{} \textbf{Cuori} e \sD{} \textbf{Quadri}: i semi rossi.\par
 \sS{} \textbf{Picche} e \sC{} \textbf{Fiori}: i semi neri.\par\medskip
Di norma il seme \textbf{non} rende una carta migliore di un’altra.
\end{minipage}
\secz{Dal più basso al più alto}
\begin{center}
 \tikz{\foreach \v [count=\i from 0] in {2,3,4,5,6,7,8,9,10,J,Q,K,A}{
   \node[draw=gold,fill=white,minimum width=7.6mm,minimum height=9mm,inner sep=0,font=\bfseries\small,text=txt] at (\i*8.3mm,0) {\v};}}
\end{center}
Le figure sono il \textbf{Fante} (J), la \textbf{Donna} (Q) e il \textbf{Re} (K). L’\textbf{Asso} (A) è speciale: di solito è la più alta, ma in certi giochi vale anche \emph{1}. Ogni scheda lo specifica.

\secz{Il vocabolario del tavolo}
\begingroup\renewcommand{\arraystretch}{1.2}
\begin{tabularx}{\linewidth}{@{}>{\bfseries}p{23mm}>{\raggedright\arraybackslash}X@{}}
 Mazziere & chi distribuisce le carte e gioca come gli altri.\\
 Banco & nel Ventuno, il giocatore contro cui giocano tutti gli altri.\\
 Mano & dalla distribuzione delle carte al risultato. Una partita può comprendere più mani.\\
 Alzare & tagliare il mazzo; lo fa chi sta a destra del mazziere.\\
 Smazzare & distribuire le carte in senso orario, una alla volta.\\
 Carte & coperte: dorso in alto, valore nascosto. Scoperte: valore visibile.\\
 Scarto & la carta che si getta via, di solito scoperta.\\
 Piatto (\emph{pot}) & le fiches al centro del tavolo, usate per pagare le vincite.\\
\end{tabularx}\endgroup
\vfill
\ettore{Mescolate sempre con calma e fate alzare il mazzo: così nessuno ha carte truccate… compresi voi.}

%% 5–6 · ACEY DEUCEY
\startpage\label{acey}
\gameheader{1}{ACEY DEUCEY}{Dentro o fuori? Tre carte e un po’ di sangue freddo}{3–8 giocatori}{10–20 minuti}{1}{enea}
\textbf{L’idea.} Il mazziere gira due carte. Tu scommetti che la terza avrà un valore \textbf{in mezzo alle due}. Con 4 e J, per esempio, vinci con 5, 6, 7, 8, 9 o 10. I semi non contano; \textbf{l’Asso è sempre alto}.
\secz{Preparazione}
Ognuno riceve \textbf{10 fiches} e ne mette \textbf{1 al centro}: questi gettoni formano il \textbf{piatto comune}. Il mazziere mescola le 52 carte senza jolly. Comincia il giocatore alla sua sinistra.
\secz{Il tuo turno, passo per passo}
\begin{enumerate}
 \item Il mazziere gira \textbf{due carte davanti a te}. Se hanno lo stesso valore o sono consecutive, come 6 e 7, le scarta e ne gira due nuove.
 \item Puoi dire \textbf{«Passo»}: non paghi nulla e tocca al prossimo. Oppure dici quanto punti, per esempio \textbf{«Punto 2»}. Tieni per ora quelle fiches davanti a te.
 \item Il mazziere gira la \textbf{terza carta}. Confrontala con le prime due e paga o incassa come negli esempi qui sotto. Poi scartate le carte e passate al prossimo giocatore.
\end{enumerate}
\begin{center}
\examplecell{36mm}{\handm{4/C/n,J/D/n,9/S/u}}{\textbf{In mezzo: vinci.}\\Prendi 2 dal piatto.}\hfill
\examplecell{36mm}{\handm{4/C/n,J/D/n,2/H/u}}{\textbf{Fuori: perdi.}\\Metti 2 nel piatto.}\hfill
\examplecell{36mm}{\handm{4/C/n,J/D/n,J/S/u}}{\textbf{Uguale: «palo».}\\Metti 4 nel piatto.}
\par\vspace{1mm}{\footnotesize Negli esempi hai puntato 2. Se vinci, tieni anche le tue 2 fiches.}
\end{center}
\begin{callout}
\textbf{Non sommare le carte}: conta solo l’ordine dei valori. Prima di giocare, leggete i limiti delle puntate nella pagina seguente.
\end{callout}
\startpage
\bigtitle{Acey Deucey, in pratica}{quanto puntare e come proseguire}
\secz{Quanto puoi puntare?}
La puntata minima è \textbf{1 fiche}. Devi rispettare due limiti:
\begin{itemize}
 \item \textbf{Il piatto deve poterti pagare.} Se contiene 5 fiches, non puoi puntarne più di 5.
 \item \textbf{Devi poter pagare anche il palo.} Per puntare 2 devi avere almeno 4 fiches: 2 per la puntata e altre 2 se la terza carta è uguale a una delle prime due.
\end{itemize}
Perciò, se hai 7 fiches, puoi puntarne al massimo \textbf{3}: ne servirebbero 8 per puntarne 4. Se hai meno di 2 fiches, puoi solo passare.
\secz{Un turno con una puntata di 2}
Hai \textbf{8 fiches} e nel piatto ce ne sono \textbf{5}. Escono 4 e J. Dici «Punto 2» e tieni le tue fiches davanti a te. Il mazziere gira la terza carta:
\begin{itemize}
 \item Esce \textbf{9}: è fra 4 e J. Prendi 2 fiches dal piatto e tieni le tue: ora ne hai \textbf{10}.
 \item Esce \textbf{2}: è fuori dall’intervallo. Metti 2 delle tue fiches nel piatto: te ne restano \textbf{6}.
 \item Esce \textbf{J}: è un palo. Metti nel piatto 4 fiches in tutto: te ne restano \textbf{4}.
\end{itemize}
\secz{Quando tocca agli altri?}
Dopo il tuo turno, le carte usate vanno negli scarti. Il mazziere ne gira due nuove per il giocatore successivo, in senso orario. Dopo un turno a testa, \textbf{mazziere compreso}, il giro è finito e il ruolo di mazziere passa a sinistra.
\secz{Fine della partita}
Dopo \textbf{tre giri}, dividete il piatto fra chi è ancora in gioco. Gli eventuali resti vanno, uno a testa, da sinistra del mazziere. Vince chi ha più fiches; a pari numero si condivide la vittoria. Se resta una sola persona, la partita finisce subito e prende il piatto.
\begin{callout}
\textbf{Piatto vuoto:} tutti aggiungono 1 fiche; chi non può esce.\\
\textbf{Mazzo quasi finito:} se restano meno di tre carte, rimescolate gli scarti prima del turno.\\
\textbf{Asso e 2:} una carta dal 3 al Re vince; un altro Asso o un altro 2 fa palo.
\end{callout}

%% 7–8 · VENTUNO
\startpage\label{ventuno}
\gameheader{2}{VENTUNO}{Blackjack: battere il banco senza sballare}{2–8 giocatori}{20–40 minuti}{1}{ettore}
\textbf{L’idea.} Avvicinarsi il più possibile a \textbf{21} senza superarlo, battendo il \emph{banco}, che è uno dei giocatori. Ci si sfida contro il banco, non tra di voi. Il banco passa a sinistra a ogni mano, così tutti hanno lo stesso vantaggio.
\secz{Quanto valgono le carte}
\begin{center}
\examplecell{34mm}{\hand{7/S,4/D}}{\textbf{2–10}\\Il loro numero}\hfill
\examplecell{38mm}{\hand{K/H,Q/C}}{\textbf{Fante, Donna, Re}\\10 punti}\hfill
\examplecell{38mm}{\hand{A/D}}{\textbf{Asso}\\1 oppure 11}
\end{center}
Il seme non conta. L’Asso vale 11 se non fa sballare, altrimenti 1: per esempio, \textbf{Asso e 6 fanno 17}; pescando un 9 diventano \textbf{16}, perché l’Asso conta 1.
\secz{Preparazione}
Ognuno parte con \textbf{20 fiches}. In ogni mano la puntata è fissa: \textbf{2 fiches} a testa. Un giocatore fa il banco; sceglietelo a sorte. Chi sta a destra del banco alza il mazzo.
\begin{callout}
\textbf{Blackjack} = Asso più 10, J, Q o K nelle prime due carte. Batte un 21 fatto con più carte e paga \textbf{3 a 2}: puntando 2, riprendi le tue 2 e ne vinci 3.
\end{callout}
\secz{Distribuzione e controllo}
\begin{enumerate}
 \item I giocatori mettono 2 fiches davanti a sé. Il banco dà \textbf{due carte scoperte} a ciascuno e due a sé: una scoperta e una coperta.
 \item Con Asso o valore da 10 scoperto, il banco controlla la carta coperta. Se ha Blackjack, lo mostra e incassa le puntate. Chi ha a sua volta Blackjack riprende le proprie 2 fiches: è un pareggio.
 \item Se il banco non ha Blackjack, paga subito gli eventuali Blackjack dei giocatori: chi è già stato pagato \textbf{non chiede altre carte e non viene pagato di nuovo}.
\end{enumerate}

\startpage
\secz{Carta o sto?}
\begin{enumerate}[resume]
 \item Da sinistra del banco, ciascuno sceglie \textbf{«Carta!»}: il banco gli aggiunge una carta scoperta. Può chiederne ancora o dire \textbf{«Sto!»}, fermandosi. Chi supera 21 perde subito la puntata, anche se poi sballa il banco.
 \item Il banco scopre la carta coperta: \textbf{pesca sotto 17 e si ferma da 17 in su}, anche con Asso e 6.
 \item Ogni giocatore rimasto confronta il proprio totale con quello del banco. Il banco incassa o paga con le proprie fiches, seguendo la tabella.
\end{enumerate}
\begin{center}\small\setlength{\tabcolsep}{2mm}\renewcommand{\arraystretch}{1.35}
\begin{tabular}{@{}m{50mm}m{62mm}@{}}
\rowcolor{tint}\textbf{Se…} & \textbf{allora…}\\
Tu sballi & perdi le 2 fiches puntate\\
Il banco sballa, tu no & riprendi 2 e vinci altre 2\\
Il tuo totale è maggiore & riprendi 2 e vinci altre 2\\
I totali sono uguali & riprendi le tue 2: pareggio\\
Il tuo totale è minore & perdi le 2 fiches puntate\\
\end{tabular}
\end{center}
\begin{center}
\examplecell{53mm}{\hand{A/S,K/D}}{\textbf{Blackjack: +3} netti.}\hfill
\examplecell{53mm}{\hand{10/H,6/C,8/D}}{\textbf{24: sballato, perdi la puntata (−2).}}
\end{center}
\secz{Raddoppio e gestione del mazzo}
Dopo le prime due carte puoi \textbf{raddoppiare a 4 fiches}: ricevi una sola carta, poi ti fermi. Vinci o perdi 4; un 21 con tre carte non è Blackjack. \textbf{Rimescolate tutte le carte tra una mano e l’altra}, mai durante una mano.
\begin{callout}
	\textbf{Perché «Black Jack»?} Il Ventuno è nato come passatempo da salotto nell’Europa del Settecento. Secondo la leggenda, il nome americano viene da un premio speciale (10 a 1!) dato a chi aveva l’Asso di Picche con un Jack nero; gli storici sono divisi, ma il nome è rimasto, il premio no.
\end{callout}
\vfill
\enea{Da \textbf{17 in su}, di solito conviene fermarsi. Con \textbf{11 o meno} chiedi carta senza paura: non puoi sballare. Tra 12 e 16, guarda la carta scoperta del banco: se è bassa (\textbf{da 2 a 6}), spesso conviene fermarsi e lasciare che sia il banco a rischiare.}

%% 9–11 · POKER
\startpage\label{poker}
\gameheader{3}{POKER}{A cinque carte, con cambio: il re degli speakeasy}{2–6 giocatori}{30–60 minuti}{2}{enea}

\textbf{L’idea.} Formare la migliore mano di \textbf{cinque carte}, con un solo cambio. Si usa il \textbf{mazzo completo da 52}, senza jolly. Per vincere puoi avere le carte migliori… oppure convincere tutti a ritirarsi.

\begin{callout}
	Negli anni Venti si giocava così (\textbf{\emph{Five-Card Draw}}): il Texas Hold’em arriverà a Las Vegas solo negli anni Sessanta!
\end{callout}
\vfill
\secz{Preparazione}
\textbf{20 fiches a testa}; il mazziere passa a sinistra a ogni mano. La prima puntata e ogni aumento sono di \textbf{1 fiche}, con al massimo \textbf{tre rilanci complessivi per giro di puntate}.
\secz{Le parole del poker}
\begingroup\renewcommand{\arraystretch}{1.15}
\begin{tabularx}{\linewidth}{@{}>{\bfseries}p{17mm}>{\raggedright\arraybackslash}X@{}}
 Check & «parola»: resti in gioco senza puntare (se nessuno ha puntato).\\
 Bet & sei il primo a puntare nel giro: metti 1 fiche nel piatto.\\
 Call & «vedo»: pareggi la puntata più alta del giro.\\
 Raise & «rilancio»: pareggi la puntata più alta, poi aggiungi ancora 1 fiche.\\
 Fold & ti ritiri e lasci nel piatto quanto hai già versato.\\
\end{tabularx}\endgroup
\secz{Una mano, passo per passo}
\begin{enumerate}
 \item \textbf{Quota iniziale.} Prima di ricevere le carte, ognuno mette 1 fiche al centro del tavolo. Poi il mazziere dà \textbf{5 carte coperte} a ciascuno: guardate solo le vostre.
 \item \textbf{Primo giro di puntate:} comincia chi è a sinistra del mazziere. Si prosegue in senso orario finché tutti hanno agito e pareggiato l’ultima puntata, oppure hanno fatto fold. Se tutti fanno check, si procede comunque.
 \item \textbf{Cambio:} da sinistra del mazziere, chi è ancora in gioco sceglie da \textbf{0 a 3 carte} da restituire coperte al mazziere e ne riceve lo stesso numero dal mazzo.
 \item \textbf{Secondo giro:} stesse regole; apre il primo giocatore ancora in gioco a sinistra del mazziere.
 \item \textbf{Showdown:} chi è rimasto mostra le carte; vince la mano migliore (classifica nella pagina seguente). Se resta un solo giocatore, incassa subito senza mostrare nulla.
\end{enumerate}
\startpage
\bigtitle{La classifica delle mani}{dalla più forte alla più debole}
\vspace{-1mm}
% Rombo di misura fissa (non dipende dalla cifra), staccato dal testo;
% descrizione con l'interlinea del suo corpo (\par dentro il gruppo);
% stesso spazio sopra e sotto il filetto (sotto si somma \parskip).
\newcommand{\rankrow}[4]{%
 \noindent\begin{minipage}[c]{9mm}\raggedright\tikz\node[diamond,draw=gold,fill=ink,inner sep=0pt,minimum size=6.4mm,text=cream,font=\decob\bfseries\small]{#1};\end{minipage}%
 \begin{minipage}[c]{\dimexpr\linewidth-9mm-41mm\relax}\raggedright{\bfseries #2\par}\vspace{.3mm}{\scriptsize #3\par}\end{minipage}\hfill
 \begin{minipage}[c]{39mm}\raggedleft\scalebox{.78}{\hand{#4}}\end{minipage}\par
 \vspace{1.5mm}{\color{gold}\hrule height .35pt}\vspace{\dimexpr1.5mm-\parskip\relax}}
\rankrow{1}{Scala colore}{Cinque carte in sequenza, tutte dello stesso seme. La più alta (10-J-Q-K-A) è la \emph{Scala reale}.}{10/H,J/H,Q/H,K/H,A/H}
\rankrow{2}{Poker}{Quattro carte dello stesso valore, più una qualsiasi.}{9/S,9/H,9/D,9/C,K/S}
\rankrow{3}{Full}{Un tris più una coppia.}{Q/S,Q/H,Q/D,7/C,7/S}
\rankrow{4}{Colore}{Cinque carte dello stesso seme, non in sequenza.}{A/C,J/C,8/C,5/C,2/C}
\rankrow{5}{Scala}{Cinque carte in sequenza, non tutte dello stesso seme. Asso basso (A-2-3-4-5) o alto (10-J-Q-K-A), ma non in mezzo.}{5/S,6/D,7/C,8/H,9/S}
\rankrow{6}{Tris}{Tre carte dello stesso valore.}{8/S,8/H,8/D,K/C,3/S}
\rankrow{7}{Doppia coppia}{Due coppie diverse, più una carta.}{J/S,J/D,4/C,4/H,A/D}
\rankrow{8}{Coppia}{Due carte dello stesso valore.}{10/H,10/C,A/S,7/D,3/C}
\rankrow{9}{Carta alta}{Nessuna combinazione: conta la carta più alta.}{A/H,J/C,8/D,6/S,3/D}
\vfill
\begin{center}\itshape\footnotesize\color{golddk}Il seme non conta mai nel confronto tra mani: a parità di combinazione si guarda il valore delle carte.\end{center}


%% 11 · SPAREGGI E CAMBIO
\startpage
\secz{Chi vince a parità di combinazione?}
Si confrontano i \textbf{valori}, partendo dalla carta più importante della combinazione:
\begin{itemize}
	\item \textbf{Coppia, tris, poker, full:} vince il valore più alto della combinazione (un tris di 9 batte un tris di 7; nel full conta il tris).
	\item \textbf{Doppia coppia:} prima la coppia più alta, poi la seconda, poi la carta rimasta.
	\item \textbf{Scala, colore, carta alta:} vince la carta più alta. La scala A-2-3-4-5 è la più bassa, la 10-J-Q-K-A la più alta.
	\item \textbf{Stessa mano identica:} il piatto si divide. I semi non servono mai a spareggiare.
\end{itemize}
\secz{Un esempio di cambio}
Ricevi queste carte. Hai due Assi: \textbf{tieni la coppia}, scarta le altre tre carte e pesca.
\begin{center}
\examplecell{55mm}{\handm{A/S/u,A/D/u,7/C/d,4/H/d,2/S/d}}{\itshape Prima: tieni i due Assi.}\hfill
\examplecell{55mm}{\handm{A/S/n,A/D/n,A/H/u,9/C/u,K/D/u}}{\itshape Dopo: un tris di Assi!}
\end{center}
\secz{Quando finisce la partita?}
Un giro di mazziere è completo quando tutti hanno distribuito una volta. Si giocano \textbf{tre giri}; chi è uscito viene saltato. Per iniziare una mano servono \textbf{almeno 9 fiches}: coprono la quota iniziale e massimo 4 fiches in ciascuno dei due giri di puntate. Chi ne ha meno esce, conservando il proprio saldo. La partita finisce prima se restano meno di due giocatori. Vince chi ha più fiches; a pari merito si condivide la vittoria.
\secz{Quanto metto per vedere?}
\begin{callout}
\textbf{Esempio di call:} hai già messo 1 fiche nel giro e il totale da vedere è salito a 3? Ne aggiungi solo \textbf{2}. La quota iniziale non conta come puntata del giro.
\end{callout}
\vfill
\ettore{Il bluff ha un prezzo: se qualcuno vede, le carte decidono. Non serve rilanciare a ogni turno; scegli quando rischiare e tieni un’espressione indecifrabile.}

%% 12–14 · GIN RUMMY
\startpage\label{gin}
\gameheader{4}{GIN RUMMY}{Il duello elegante: pescare, scartare, combinare}{2 giocatori}{30–45 minuti}{3}{ettore}
\textbf{L’idea.} Con \textbf{10 carte}, formare tris e scale e ridurre i punti delle carte \textbf{spaiate}, cioè fuori dalle combinazioni. Quando ne restano pochi, puoi «bussare» e chiudere la mano. Si gioca \textbf{in due}, fino a 100 punti.

\secz{Preparazione}
Un mazzo da \textbf{52 carte senza jolly}. Si gioca in due. Scegliete a sorte il primo mazziere, poi alternatevi. \textbf{Anche il mazziere gioca}: non c’è un banco.
\secz{Quanto valgono le carte}
\begin{itemize}
 \item \textbf{Asso = 1}, sempre basso.
 \item \textbf{Dal 2 al 10}: il loro numero.
 \item \textbf{Fante, Donna e Re = 10}.
\end{itemize}
\secz{Le combinazioni}
\begin{itemize}
 \item \textbf{Tris o poker:} 3 o 4 carte dello stesso valore (come \cardref{7}{\sS}, \cardref{7}{\sH}, \cardref{7}{\sD}).
 \item \textbf{Scala:} almeno 3 carte dello stesso seme in sequenza (come \cardref{9}{\sC}, \cardref{10}{\sC}, \cardref{J}{\sC}, \cardref{Q}{\sC}). L’Asso sta solo \textbf{all’inizio}: A-2-3 vale, Q-K-A no.
\end{itemize}
Una carta può appartenere a \textbf{una sola combinazione}. Il totale delle spaiate serve a capire quando puoi chiudere: \textbf{non va aggiunto al tuo punteggio di partita}. I punti che guadagni si calcolano a pagina \pageref{gin-esempio}.
\secz{Distribuzione e prima mossa}
Il mazziere dà \textbf{10 carte a testa}, ne gira una per iniziare gli \textbf{scarti} e lascia il resto coperto come \textbf{mazzo di pesca}.

Chi non ha distribuito può prendere la carta scoperta o rifiutarla. Se rifiuta, può prenderla il mazziere. Chi la prende \textbf{scarta e passa il turno}. Se entrambi rifiutano, chi non ha distribuito pesca dal mazzo e scarta: poi si alternano i turni normali.

\startpage
\secz{Un turno normale}
\begin{enumerate}
 \item \textbf{Pesca}: aggiungi alla tua mano una carta, scegliendo fra la prima del mazzo coperto e quella in cima agli scarti.
 \item \textbf{Scarta}: scegli una delle tue carte e mettila scoperta sopra gli scarti. Se hai preso dagli scarti, non puoi rimettere subito la stessa carta.
\end{enumerate}
Dopo aver pescato hai 11 carte; dopo lo scarto torni a \textbf{10}. Non posare le combinazioni sul tavolo finché qualcuno chiude.
\label{gin-esempio}
\secz{Bussare («knock») e Gin}
Se, dopo lo scarto, le tue spaiate valgono \textbf{10 punti o meno}, puoi dichiarare «Busso!» e mostrare la mano. Con \textbf{zero punti spaiati} fai \textbf{Gin}. Non sei obbligato a bussare appena puoi.
\begin{center}
\examplecell{32mm}{\hand{7/S,7/H,7/D}}{\textbf{Tris}}\hfill
\examplecell{40mm}{\hand{9/C,10/C,J/C,Q/C}}{\textbf{Scala}}\hfill
\examplecell{40mm}{\hand{A/D,3/H,5/S}}{\textbf{Spaiate: 1\,+\,3\,+\,5\,=\,9}}
\end{center}
Con questa mano puoi bussare: le tre spaiate valgono 9.
\secz{L’attacco dopo la bussata}
Dopo che hai bussato, l’avversario mostra la mano e forma le sue combinazioni. Può poi \textbf{aggiungere le sue spaiate alle tue combinazioni}: per esempio un K\sC{} alla scala qui sopra, o un 7\sC{} al tris. Questo si chiama «attaccare» e riduce le sue spaiate. Tu \textbf{non puoi aggiungere carte alle sue combinazioni}. Con il Gin, invece, \textbf{nessuno può attaccare}.
\secz{Quanti punti si fanno}
In questa versione, Gin e controbussata hanno entrambi \textbf{25 punti di bonus}. Segna solo chi vince la mano:
\begin{itemize}
 \item \textbf{Bussata vincente:} se, dopo gli attacchi, hai meno punti spaiati dell’altro, segni la \textbf{differenza}.
 \item \textbf{Gin:} segni le spaiate avversarie \textbf{più 25}; vinci anche se l’altro ha zero spaiate.
 \item \textbf{Controbussata (\emph{undercut}):} senza Gin, se l’altro resta con spaiate \textbf{uguali o inferiori} alle tue, segna lui la differenza \textbf{più 25}.
\end{itemize}

\startpage
\secz{Fine della mano e della partita}
Se chi pesca lascia \textbf{due carte nel mazzo} e poi scarta senza chiudere, la mano è nulla: nessun punto, rimescolate e ridistribuisce lo stesso mazziere. Altrimenti, il ruolo cambia a ogni mano. Vince la partita chi raggiunge o supera \textbf{100 punti}.
\secz{Esempi di punteggio}
\begin{callout}
Bussi con \textbf{9}, l’altro resta a \textbf{14}: segni 14\,−\,9\,=\,5. Se invece scende a \textbf{6}, segna lui 9\,−\,6\,+\,25\,=\,28. A \textbf{9 contro 9}, l’altro segna 25.
\end{callout}
\secz{Un attacco che cambia il risultato}
Nell’esempio di pagina \pageref{gin-esempio} bussi con 9. Dopo le proprie combinazioni, all’altro restano \cardref{7}{\sC}, \cardref{K}{\sC}, \cardref{2}{\sH} e \cardref{2}{\sC}: 21 punti. Attacca il 7 al tuo tris e il Re alla tua scala: resta con soli \textbf{4 punti}. Ti controbussa e segna 9\,−\,4\,+\,25\,=\,30. Ecco perché bussare comporta un rischio!

\secz{Un po’ di strategia}
\begin{itemize}
 \item \textbf{Scarta le figure isolate}: valgono 10 punti ciascuna.
 \item Conserva coppie e carte vicine dello stesso seme: possono diventare combinazioni.
 \item Guarda cosa raccoglie l’altro dagli scarti: una carta vicina potrebbe aiutarlo.
 \item Non bussare solo perché puoi: se l’altro ha poche spaiate, rischi la controbussata.
\end{itemize}
\secz{Riepilogo veloce}
\begin{callout}
\begin{enumerate}
 \item Dieci carte a testa. Pesca, scarta e tieni le combinazioni in mano.
 \item Con al massimo 10 punti spaiati puoi bussare.
 \item L’altro può attaccare: confrontate le spaiate \textbf{dopo} gli attacchi.
 \item Il Gin blocca gli attacchi e vince sempre; la controbussata premia chi non ha bussato. Entrambi danno 25 di bonus.
 \item Sommate le mani fino a 100 punti.
\end{enumerate}
\end{callout}
\secz{Un po’ di storia}
Il Gin Rummy discende dal Rummy ed è stato inventato nel 1909 da Elwood T. Baker. È rimasto uno dei giochi per due più amati: partite rapide, regole chiare e tanta tattica. Ideale per sfidare la persona accanto a voi, tra un ballo e l’altro.
\vfill
\ettore{Bussare è una scelta di tempo: chiudi prima che l’altro migliori, ma gli concedi un ultimo attacco. Il Gin richiede più pazienza e toglie questa possibilità all’avversario.}

%% 15–17 · CUORI
\startpage\label{cuori}
\gameheader{5}{CUORI}{Hearts: l’arte di non prendere punti}{4 giocatori (o 3)}{60–90 minuti}{3}{enea}
\textbf{L’idea.} Qui si vince facendo il \textbf{minor numero di punti}. Ogni \textbf{Cuore} che prendi vale \textbf{1 punto}, la \textbf{Donna di Picche} ne vale \textbf{13}: gli altri semi non valgono nulla. È un gioco di «prese»: l’obiettivo è schivarle con stile. Ogni mano vale 26 punti in tutto, da distribuire (a malincuore) tra i giocatori.

\begin{callout}
\textbf{Che cos’è una presa?} Tutti giocano una carta a testa. Chi vince raccoglie tutte le carte giocate: quelle con Cuori o Donna di Picche gli danno penalità. «Aprire» significa giocare la prima carta della presa.
\end{callout}
\secz{Preparazione}
In \textbf{quattro}, ognuno per sé, con 52 carte senza jolly. Il mazziere dà \textbf{13 carte a testa} e passa a sinistra a ogni mano. Il 2 è basso, l’Asso alto. \textbf{Non c’è briscola}: vince la carta più alta del seme con cui si apre la presa.
\secz{Il passaggio delle carte}
Ognuno sceglie \textbf{3 carte della propria mano} e le dà coperte al giocatore indicato qui sotto. Riceve altre tre carte: può guardarle solo dopo aver consegnato le proprie. La direzione cambia a ogni mano:
\begin{enumerate}
	\item mano 1: si passa a \textbf{sinistra};
	\item mano 2: a \textbf{destra};
	\item mano 3: \textbf{di fronte};
	\item mano 4: \textbf{nessun passaggio}, poi si ricomincia.
\end{enumerate}

\secz{La prima presa}
Chi ha il \textbf{2 di Fiori} deve aprire con quella carta. Gli altri devono giocare una carta di Fiori, se ne hanno: si dice \textbf{«rispondere al seme»}. Chi non ha Fiori gioca una carta diversa, ma \textbf{non un Cuore né la Donna di Picche}, a meno che gli restino soltanto carte di questi due tipi.

\startpage
\secz{Le prese successive}
\begin{itemize}
 \item Chi ha vinto la presa precedente \textbf{apre} la nuova.
 \item Devi \textbf{rispondere al seme}, se puoi. Solo se non hai nessuna carta di quel seme puoi giocarne una diversa, anche con penalità: non potrà vincere quella presa.
 \item La carta più alta del \textbf{seme di apertura} vince: chi la gioca raccoglie tutte le carte della presa, con i relativi punti.
\end{itemize}

\begin{center}
\examplecell{30mm}{\handm{8/C/n}}{\itshape Apre}\hfill
\examplecell{30mm}{\handm{3/C/n}}{\itshape Risponde}\hfill
\examplecell{30mm}{\handm{K/C/u}}{\itshape Vince}\hfill
\examplecell{30mm}{\handm{5/H/n}}{\itshape Non ha Fiori}
\end{center}

\textbf{Esempio dopo la prima presa.} Il K\sC{} è il Fiori più alto \textbf{tra le carte giocate}: prende anche il 5\sH{} e quindi 1 punto. Un Asso di Fiori avrebbe battuto il Re.
\secz{Cuori «rotti» e Donna di Picche}
\begin{itemize}
 \item Non puoi \textbf{aprire con un Cuore} finché qualcuno non ha giocato un Cuore in una presa precedente, salvo che tu abbia solo Cuori. Da quel momento i Cuori sono «rotti».
 \item In questa versione, la sola \textbf{Donna di Picche non rompe i Cuori}.
 \item Dopo la prima presa, puoi aprire con la Donna di Picche; per scartarla su un altro seme devi però \textbf{non poter rispondere} a quel seme.
\end{itemize}
\secz{Finita la mano}
Dopo le 13 prese, ognuno conta i punti raccolti: \textbf{1 per ogni Cuore}, \textbf{13 per la Donna di Picche}. Controllo: i punti di tutti devono fare \textbf{26} in totale.

\secz{La Luna (\emph{shooting the moon})}
Se un giocatore prende \textbf{tutti i Cuori e la Donna di Picche}, sbanca la mano: lui fa \textbf{0 punti} e ciascuno degli altri tre fa \textbf{26 punti}.

\startpage
\secz{Fine della partita}
Alla fine di ogni mano aggiungete le penalità ai totali precedenti. Se qualcuno raggiunge o supera \textbf{100 punti}, la partita finisce: vince chi ha \textbf{meno punti}. In caso di pari punteggio minimo, la vittoria è condivisa.
\secz{Esempio di punteggio}
\begin{center}\small\setlength{\tabcolsep}{2.5mm}\renewcommand{\arraystretch}{1.45}
\begin{tabular}{@{}lrrrr@{}}
\rowcolor{tint}& \textbf{Anna} & \textbf{Bruno} & \textbf{Carla} & \textbf{Dario}\\
Mano 1 & 4 & 13 & 8 & 1\\
Mano 2 & 0 & 7 & 6 & 13\\
Mano 3 & 12 & 0 & 14 & 0\\\hline
\textbf{Totale} & \textbf{16} & \textbf{20} & \textbf{28} & \textbf{14}\\
\end{tabular}
\end{center}
Qui non ci sono Lune: ogni mano somma 26. Dario è in testa con 14.
\secz{In tre?}
\begin{itemize}
 \item Togliete il \textbf{2 di Quadri}: restano 51 carte. Distribuitene \textbf{17 a testa} e giocate \textbf{17 prese}.
 \item Passate tre carte a \textbf{sinistra}, poi a \textbf{destra}, poi \textbf{nessuna}; ripetete questo ciclo di tre mani.
 \item Parte sempre il \textbf{2 di Fiori}. Le altre regole non cambiano.
 \item Con la Luna, chi la fa segna 0 e gli altri \textbf{due} segnano 26: totale 52. Senza Luna il totale resta 26.
\end{itemize}
\secz{Consigli per cominciare}
\begin{itemize}
 \item Nel passaggio, cerca di \textbf{svuotare un seme}: potrai scartare penalità quando quel seme viene aperto, salvo i vincoli della prima presa.
 \item Tieni qualche carta bassa per evitare le prese pericolose. Non tutte le prese fanno male: quelle senza penalità valgono zero.
 \item Segui la Donna di Picche e guarda se qualcuno raccoglie tutti i Cuori: potrebbe tentare la Luna.
\end{itemize}
\vfill
\ettore{Non basta perdere una presa: bisogna perdere quella giusta. Prima di giocare una carta alta, guarda quanti Cuori sono già sul tavolo.}

%% 18–19 · BONUS RAMINO
\startpage\label{ramino}
\gameheader{BONUS}{RAMINO}{Una variante semplice, con due mazzi e quattro jolly}{2–5 giocatori}{30–45 minuti}{2}{ettore}
\begin{callout}
{\suitfont\color{golddk}★}\quad\textbf{Gioco bonus: servono due mazzi!} Chiedete a un altro tavolo di unire le forze: mescolate insieme i due mazzi, \emph{compresi i 4 jolly}. Più mazzi, più divertimento..
\end{callout}
\textbf{L’idea.} Il Ramino è il capostipite dei giochi di «combinazioni»: già diffuso nei primi del Novecento, da lui discende il Gin Rummy che avete appena imparato. Il gioco è semplice: formare tris e scale e liberarsi per primi di tutte le carte. Questa è la nostra \textbf{variante semplificata}.
\secz{Preparazione}
\textbf{108 carte}: due mazzi da 52 più quattro jolly. Da 2 a 5 giocatori, \textbf{13 carte a testa}, da tenere nascoste agli altri. Il mazziere viene scelto a sorte e cambia a ogni mano, passando a sinistra. Comincia chi è alla sua sinistra; si procede in senso orario.

Girate una carta per iniziare gli \textbf{scarti}. Se è un jolly, cercate una carta diversa e poi rimescolate i jolly nel mazzo coperto rimasto: sarà il \textbf{mazzo di pesca}.
\secz{Le combinazioni}
\begin{itemize}
 \item \textbf{Tris o poker:} 3 o 4 carte dello stesso valore, di \textbf{semi distinti}. Due 9 di Picche non possono stare nello stesso gruppo.
 \item \textbf{Scala:} almeno 3 carte dello stesso seme in sequenza. Asso all’inizio (A-2-3) o alla fine (Q-K-A); K-A-2 non vale.
 \item \textbf{Jolly:} sostituisce una carta; massimo \textbf{uno per combinazione}. Dichiarate quale carta rappresenta.
\end{itemize}
\begin{center}
\examplecell{52mm}{\hand{9/S,9/H,9/D}}{\textbf{Tris di 9}, semi distinti}\hfill
\examplecell{58mm}{\hand{4/C,5/C,JK/Z,7/C}}{\textbf{Scala}: il jolly è il 6\sC{}}
\end{center}

\startpage
\secz{Un turno}
\begin{enumerate}
 \item \textbf{Pesca}: aggiungi alla tua mano una carta, scegliendo fra la prima del mazzo coperto e quella in cima agli scarti.
 \item \textbf{Cala}, se vuoi: posa davanti a te un tris, un poker o una scala, così tutti li vedono. Puoi posarne più di uno. \textbf{Non serve un minimo di punti}.
 \item Se hai già calato, anche in questo turno, puoi \textbf{attaccare} carte alle combinazioni tue o altrui, mantenendole valide: un 6\sC{} alla scala 7-8-9\sC{}, o il quarto seme a un tris. I gruppi non superano quattro carte.
 \item \textbf{Scarta}: metti una delle tue carte sopra gli scarti, \textbf{se non hai già chiuso}; non quella appena presa dagli scarti. Poi tocca al giocatore alla tua sinistra.
\end{enumerate}
\secz{Chiudere e contare}
Chi resta senza carte, calandole, attaccandole o scartando l’ultima, \textbf{chiude e segna 0}. In questa variante \textbf{si può chiudere senza scarto}. Gli altri contano solo le carte rimaste in mano:
\begin{center}\small\setlength{\tabcolsep}{3mm}\renewcommand{\arraystretch}{1.3}
\begin{tabular}{@{}m{45mm}C{25mm}@{}}
\rowcolor{tint}\textbf{Carta} & \textbf{Punti}\\
Dal 2 al 10 & il loro numero\\
Fante, Donna, Re & 10\\
Asso & 11\\
Jolly & 20\\
\end{tabular}
\end{center}
Giocate \textbf{tre mani} e sommate le penalità: vince il totale \textbf{più basso}. In caso di parità, vittoria condivisa.
\secz{Regole della nostra variante}
\begin{itemize}
 \item I jolly già calati \textbf{non si recuperano}; non si smontano né si riordinano le combinazioni sul tavolo.
 \item Se il mazzo finisce, rimescolate gli scarti tranne l’ultimo. Se non ci sono carte da rimescolare, chi è di turno può prendere l’ultimo scarto; se non vuole, la mano finisce e \textbf{tutti contano le carte rimaste}, senza un vincitore di mano.
\end{itemize}
\vfill
\enea{Se giocate in quattro o cinque, cercate di \textbf{calare presto}: vi dà la possibilità di attaccare le carte sul tavolo e scaricare la mano. Se aspettate troppo, rischiate di restare con tanti punti in mano.}

%% 20 · GALATEO E SEGNAPUNTI
\startpage\label{segnapunti}
\bigtitle{Il galateo del tavolo}{poche regole d’oro per giocare da signori}
\begin{itemize}
 \item \textbf{Il mazziere mescola;} chi sta alla sua destra alza.
 \item \textbf{Mai sbirciare} le carte altrui. Neppure «per sbaglio».
 \item \textbf{Si vince con modestia e si perde con eleganza.}
 \item Nel Ventuno \textbf{il banco ruota}; negli altri giochi seguite la rotazione del mazziere indicata nella scheda.
 \item \textbf{Niente soldi}: solo gloria, fiches e un brindisi.
\end{itemize}
\secz{Per i più coraggiosi: il Bridge}
Nel \textbf{1925}, durante una crociera, Harold Stirling Vanderbilt mise a punto il \emph{Contract Bridge}. Si gioca in quattro, a coppie: nella licitazione i compagni descrivono la propria mano con \textbf{dichiarazioni regolamentate e accordi noti anche agli avversari}. Divenne presto una mania mondiale. Le regole? Se vi sentite temerari, cercatele a fine serata!
\secz{Segnapunti}
\begin{center}\small\arrayrulecolor{gold}\setlength{\tabcolsep}{0pt}
\renewcommand{\arraystretch}{2.05}
\begin{tabular}{|C{18mm}|C{24mm}|C{24mm}|C{24mm}|C{24mm}|}
\hline
\rowcolor{tint}\textbf{Mano} & \textbf{Nome} & \textbf{Nome} & \textbf{Nome} & \textbf{Nome}\\\hline
1 & & & &\\\hline
2 & & & &\\\hline
3 & & & &\\\hline
4 & & & &\\\hline
5 & & & &\\\hline
6 & & & &\\\hline
7 & & & &\\\hline
\rowcolor{tint}\textbf{Totale} & & & &\\\hline
\end{tabular}
\end{center}
\vfill
\begin{center}\itshape\small\color{golddk}La serata finisce, la voglia di giocare no: ai prossimi vent’anni!\end{center}

%% ------------------------------------------------------------
%% 24 · RETROCOPERTINA
%% ------------------------------------------------------------
\startpage[dark]
\begin{tikzpicture}[remember picture,overlay]
 \node at ([yshift=-72mm]current page.north) {\includegraphics[width=104mm]{img/ventanni_F.png}};
 \node[text=cream,font=\itshape\fontsize{11}{15}\selectfont,align=center,text width=104mm] at ([yshift=-137mm]current page.north)
   {Grazie di aver giocato con noi.\\Buona fortuna — o meglio, buona mano!};
 \draw[gold,line width=.4pt] ([xshift=-30mm,yshift=-153mm]current page.north) -- ([xshift=-3mm,yshift=-153mm]current page.north);
 \draw[gold,line width=.4pt] ([xshift=3mm,yshift=-153mm]current page.north) -- ([xshift=30mm,yshift=-153mm]current page.north);
 \node[diamond,fill=gold,inner sep=1.3pt] at ([yshift=-153mm]current page.north) {};
 \node[text=gold,font=\decob\fontsize{9}{13}\selectfont,align=center] at ([yshift=-168mm]current page.north)
   {SABATO 21 NOVEMBRE · 24ORE, MODENA\\[1mm]20 ANNI DI ESN MORE · GALA ANNI ’20};
 \node[text=cream!70!ink,font=\footnotesize,align=center] at ([yshift=22mm]current page.south)
   {Erasmus Student Network · Modena e Reggio Emilia\\\scriptsize Fate il vostro gioco... e che vinca il migliore.};
\end{tikzpicture}
\end{document}
